% AN OUTPUT ROUTINE CLEARS THE PARAGRAPH SHAPE INSIDE ITS OWN GROUP.
%
% A paragraph may begin inside \output, so the reference clears the four
% things that shape one -- \looseness, \hangindent, \hangafter and
% \parshape -- when it fires the routine up. It does that AFTER opening
% the routine's group, so the clearing is local: what the document set
% comes back when the routine ends, and \tracingrestores draws all four:
%
%   {restoring \parshape=3}
%   {restoring \hangafter=-2}
%   {restoring \hangindent=5.0pt}
%   {restoring \looseness=1}
%
% in that order, being the reverse of the order they were cleared in.
% Only the ones that were not already clear are saved, on the usual rule
% that an assignment which changes nothing saves nothing: with all four
% already at their defaults the routine's group restores none of them.
%
% The routine sees them cleared, not as the document left them.
\catcode`\{=1 \catcode`\}=2 \catcode`\#=6
\tracingonline=1 \tracingrestores=1
\hsize=100pt \vsize=20pt \parindent=0pt \maxdepth=2pt
\output={\message{[in: looseness=\the\looseness,hangindent=\the\hangindent,%
  hangafter=\the\hangafter,parshape=\the\parshape]}\shipout\box255 }
\looseness=1 \hangindent=5pt \hangafter=-2 \parshape=3 0pt 10pt 0pt 20pt 0pt 30pt
\message{[A a page with the four set]}
\vbox{}\vskip30pt \penalty-10000
\message{[B the four again, and clear]}
\message{[looseness=\the\looseness,hangindent=\the\hangindent,%
  hangafter=\the\hangafter,parshape=\the\parshape]}
\looseness=0 \hangindent=0pt \hangafter=1 \parshape=0
\message{[C a page with none of them set]}
\vbox{}\vskip30pt \penalty-10000
\message{[done]}
\end
